\section{Instruction flags and argument boundaries}

\subsection{\code{-I}: input and control references}

\code{-I} begins the instruction's input list:
\begin{lstlisting}
H     -I Target       -O Target
CNOT -I Control      -O Target
CCNOT -I A B          -O Target
\end{lstlisting}
Its interpretation depends on the gate. For a one-qubit gate, the input names
the same wire that the output mutates. For CNOT/CZ/CY it names one control. For
CCNOT and the two-input derived logic gates it names two controls. Input order
does not matter to symmetric predicates such as AND or XOR, but it remains part
of the written instruction and should be kept consistent for readability.

\subsection{\code{-O}: output and target references}

\code{-O} begins the output list. For most gates its first token is the target
wire whose amplitudes are changed:
\begin{lstlisting}
X -I Q -O Q
AND -I A B -O Result
\end{lstlisting}
\code{-O} does not mean ``allocate a fresh classical output.'' The target may
already contain a quantum state, and reversible logic XORs its result into that
existing state. SWAP is exceptional: it requires two outputs matching its two
inputs. MEASURE takes no output flag because the measurement map stores its
classical result.

\subsection{\code{-IF}: classical feed-forward (optional)}

A gate may carry a forward-only classical condition after a prior MEASURE:
\begin{lstlisting}
X -I Bob -O Bob -IF Alice=1
Z -I Bob -O Bob -IF M0=1
\end{lstlisting}
The named token must already have been measured; running a conditioned gate
before its MEASURE is a runtime error. If the classical bit does not match,
the gate is skipped and execution continues. The skipped gate stays in the
circuit. To condition several gates, or to pair a gate with an alternative,
use the structured \code{IF} / \code{ELSE} / \code{ENDIF} block described in
\hyperlink{if-else-blocks}{Structured classical branches}; it lowers to the
same \code{-IF} conditions. Either way there is no backward edge---the circuit
remains a finite ordered list.

\subsection{Where a flag list ends}

An \code{-I} or \code{-O} list continues until the next token beginning with a
hyphen or the end of the line:
\begin{lstlisting}
CCNOT -I A B -O Carry -$R
\end{lstlisting}
Here A and B belong to \code{-I}, Carry belongs to \code{-O}, and
\code{-$R} starts a separate flag. This rule is why token names should not
begin with a hyphen.

Flags are case-insensitive, so \code{-i} and \code{-o} work, but uppercase is
the documented style. Supplying flags in an unusual order may still parse, but
\code{-I} followed by \code{-O} is the clearest and matches generated source.

\subsection{No-parameter-substitution marker}

\begin{lstlisting}
PHASE=pi/4 -I Q -O Q -$R
\end{lstlisting}
The language recognizes this marker and, on that instruction only, resolves
names as locals instead of substituting \code{PARAMS} bindings. Aliases created
by \code{SET} still apply.

\subsection{Positional arguments}

Not every instruction uses flags. SET, MAIN-PROCESS, RETURNVALS, ACCEPTVALS,
and several compatibility instructions consume positional words:
\begin{lstlisting}
SET Target 0p
MAIN-PROCESS Example
RETURNVALS Carry Sum
\end{lstlisting}
The command name is the first word and positional arguments follow in order.
RETURNVALS order is part of the process interface. For non-gate compatibility
instructions, extra words may be retained even when no execution behavior is
defined; accepted text should not be mistaken for an implemented effect.

